%
% conversation-mode-call-agent.tex
%
% Z Specification for the persistent call-agent process lifecycle -- the
% SessionAttach implementation DES-066 proposes to replace DES-064/DES-065's
% per-turn claude -p --resume subprocess spawn.
%

\documentclass[a4paper,10pt,fleqn]{article}
\usepackage[margin=1in]{geometry}
\usepackage{fuzz}
\usepackage[colorlinks=true,linkcolor=blue,citecolor=blue,urlcolor=blue]{hyperref}

\begin{document}

\title{Vox Conversation Mode: A Z Specification for the Call-Agent Process Lifecycle}
\author{Formal Model of dormant/ready/busy/ended}
\date{August 2026}
\hypersetup{
  pdftitle={Vox Conversation Mode Call-Agent Process Lifecycle: A Z Specification},
  pdfauthor={Punt Labs},
  pdfsubject={Z Specification},
  pdfcreator={fuzz/probcli}
}
\maketitle

\tableofcontents
\newpage

\section{Introduction}

\texttt{DESIGN.md}'s DES-065 records that the mechanism behind
\texttt{SessionAttach} (a fresh \texttt{claude -p --resume} subprocess spawned
per human turn) is confirmed unworkable: 13--25s median per-turn latency,
dominated by re-paying Claude Code's plugin/hook bootstrap on every single
exchange. DES-066 proposes the replacement, and its own 2026-08-23 spike
confirmed the core mechanism live: one \texttt{pi --mode rpc} process, spawned
once per \emph{call} and held alive for its whole duration, accepting one
\texttt{prompt} command per turn over its stdin and streaming the reply back
as JSON-line events on stdout.

This document models that process's lifecycle formally, per
\texttt{docs/WORKFLOW.md}'s stateful-audio gate: the process has four states
and transitions between them, an invariant must hold across those transitions
(a turn is never dispatched while another is still in flight -- exactly the
rejection the spike observed live: \texttt{"Agent is already processing.
Specify streamingBehavior (\-`steer\-' or \-`followUp\-') to queue the
message."}), and getting a transition wrong here is exactly the class of
defect \texttt{docs/WORKFLOW.md} requires this gate to catch before
implementation, not after.

This model is deliberately narrower than, and separate from,
\texttt{docs/conversation-mode-call-state.tex}'s \texttt{CallState} machine
(\texttt{idle}/\texttt{listening}/\texttt{waiting}/\texttt{speaking}). That
model governs the call's own audio lifecycle and is unaffected by this
change -- DES-065 names it retained. This model governs a distinct object:
the single long-lived agent \emph{process} \texttt{CallState}'s
\texttt{waiting}/\texttt{speaking} transitions dispatch into and read replies
from, via \texttt{SessionAttach.send\_turn}. \S\ref{sec:relation} states the
relationship between the two explicitly, rather than merging them into one
schema, because they have different lifetimes: one \texttt{CallAgentState}
instance persists for an entire call; \texttt{CallState}'s \texttt{mode}
cycles through \texttt{listening}/\texttt{waiting}/\texttt{speaking} many
times within that same call.

\section{Basic Types}

\begin{zed}
  [TURN]
\end{zed}

$TURN$ is the same opaque handle \texttt{conversation-mode-call-state.tex}
uses -- a transcribed human utterance. This model does not re-derive it; it
only needs to know a turn is dispatched and, later, that its reply completed.

\section{Free Types}

\begin{zed}
  AgentProc ::= dormant | ready | busy | ended
\end{zed}

\begin{zed}
  ZBOOL ::= ztrue | zfalse
\end{zed}

Four states, directly evidenced by the spike's own RPC event stream:
\texttt{dormant} (no process spawned yet), \texttt{ready} (process alive,
context delivered, waiting for a prompt -- observed as the state after
\texttt{agent\_settled}), \texttt{busy} (a prompt is in flight -- observed as
the state between a \texttt{prompt} command's \texttt{response} and its
matching \texttt{agent\_end}), and \texttt{ended} (process torn down, terminal).

\section{State}
\label{sec:state}

\begin{schema}{CallAgentState}
  procState : AgentProc \\
  contextDelivered : ZBOOL \\
  readOnly : ZBOOL \\
  turnsHandled : \nat
  \where
  procState \in \{ready, busy\} \implies contextDelivered = ztrue \\
  procState \neq dormant \implies readOnly = ztrue \\
  procState = dormant \implies turnsHandled = 0
\end{schema}

Three invariants, each pinning an open question DES-066 named:

\begin{enumerate}
  \item \textbf{Context is never absent while the process can accept or is
    processing a turn.} \texttt{contextDelivered} must be \texttt{ztrue}
    whenever \texttt{procState} is \texttt{ready} or \texttt{busy} -- there is
    no reachable state where a turn could be sent to a process that never
    received its opening context snapshot (DES-066 \S``Proposed decision''
    item 3).
  \item \textbf{The process is read-only against the codebase for its entire
    live extent, not just during a turn.} \texttt{readOnly} is \texttt{ztrue}
    for every state except \texttt{dormant} -- the restriction does not
    toggle on and off per turn; it is a property of the process having been
    spawned as a call agent at all (DES-066 \S``Proposed decision'' item 4).
    This model does not choose \emph{how} read-only is enforced (DES-066
    still lists that as an open unknown) -- it only fixes \emph{when} it must
    hold, so an implementation cannot satisfy the model while enforcing
    read-only only during \texttt{busy} and leaving a window open in
    \texttt{ready}.
  \item \textbf{A dormant process has handled no turns.} Trivial, but it
    blocks a spurious initial state where \texttt{turnsHandled} starts
    nonzero -- the same discipline
    \texttt{conversation-mode-call-state.tex}'s \texttt{hasPendingAddendum}
    invariant applies to \texttt{mode = idle}.
\end{enumerate}

\section{Initialisation}

\begin{schema}{Init}
  CallAgentState'
  \where
  procState' = dormant \\
  contextDelivered' = zfalse \\
  readOnly' = zfalse \\
  turnsHandled' = 0
\end{schema}

\section{Operations}

\subsection{Process lifecycle}

\begin{schema}{Spawn}
  \Delta CallAgentState
  \where
  procState = dormant \\
  procState' = ready \\
  contextDelivered' = ztrue \\
  readOnly' = ztrue \\
  turnsHandled' = 0
\end{schema}

One process is spawned per call (DES-066's central claim, superseding
DES-064/DES-065's per-turn spawn) and immediately receives its opening
context snapshot -- \texttt{Spawn} does not leave an intermediate state where
the process exists but has no context; \texttt{contextDelivered'} is set in
the same operation. \texttt{readOnly'} is set here too, per
\S\ref{sec:state}'s second invariant: read-only takes effect at spawn, not
at the first turn.

\begin{schema}{SendTurn}
  \Delta CallAgentState \\
  turn? : TURN
  \where
  procState = ready \\
  procState' = busy
\end{schema}

The precondition \texttt{procState = ready} is not a design choice this model
invents -- it is the exact behavior the 2026-08-23 spike observed live: a
\texttt{prompt} command sent while \texttt{procState} is already \texttt{busy}
is rejected by \texttt{pi}'s own RPC mode
(\texttt{"Agent is already processing"}), not queued or silently dropped. A
correct \texttt{SessionAttach} implementation must therefore never call
\texttt{send\_turn} while a prior call's reply stream has not yet completed --
this precondition is the formal statement of that caller obligation, the same
role \texttt{docs/conversation-mode-call-state.tex}'s single-serialized-dispatch-point
discipline (\S\ref{sec:relation}) already plays for \texttt{CallState}.

\begin{schema}{ReplyComplete}
  \Delta CallAgentState
  \where
  procState = busy \\
  procState' = ready \\
  turnsHandled' = turnsHandled + 1
\end{schema}

The spike's own \texttt{agent\_end}/\texttt{agent\_settled} event pair marks
this transition: the assistant's reply finished, and the process returns to
\texttt{ready} for the next turn without a fresh spawn -- this is the
persistence property the spike's \texttt{cacheRead} evidence
(turn two reused turn one's warm context) confirms actually happens, not just
that the process theoretically could.

\begin{schema}{Teardown}
  \Delta CallAgentState
  \where
  procState \neq busy \\
  procState \neq dormant \\
  procState' = ended
\end{schema}

Teardown is safe from \texttt{ready} (the ordinary end-of-call path) but
\emph{not} from \texttt{busy} -- the precondition \texttt{procState \neq busy}
forbids killing the process while a reply is still streaming, per the same
"a pending human utterance is never silently dropped" discipline
\texttt{conversation-mode-call-state.tex}'s \texttt{hasPendingAddendum}
invariant applies on the audio side. A call that hangs up mid-reply must
drain or abort the in-flight turn first (\texttt{AbortTurn}, below), not
tear the process down out from under it.

\begin{schema}{AbortTurn}
  \Delta CallAgentState
  \where
  procState = busy \\
  procState' = ended
\end{schema}

The one path that \emph{does} end the process from \texttt{busy}: an explicit
abort (the call hung up, or FR-2's inactivity timeout fired) while a reply was
still in flight. This is a distinct operation from \texttt{Teardown}, not a
relaxed precondition on it, because the two have different real-world
consequences the implementation must account for differently -- an ordinary
\texttt{Teardown} discards nothing in flight; an \texttt{AbortTurn} discards a
partial reply the human never heard, which \texttt{CallSession} (the
\texttt{CallState}-side caller, \S\ref{sec:relation}) must handle explicitly
rather than silently swallow.

\subsection{Observation}

\begin{schema}{CanAcceptTurn}
  \Xi CallAgentState \\
  canAccept! : ZBOOL
  \where
  canAccept! = ztrue \iff procState = ready
\end{schema}

A non-mutating query a caller can use to decide whether \texttt{SendTurn}'s
precondition currently holds, without attempting the call speculatively and
handling the rejection reactively -- the same
\texttt{"already processing"} case the spike hit when it fired a second
prompt before the first one's turn had completed.

\section{Relationship to the Call's Own Audio State Machine}
\label{sec:relation}

\texttt{CallSession} owns one \texttt{CallState} (audio mode) and one
\texttt{CallAgentState} (this model) for the duration of a call, with a fixed
correspondence between them:

\begin{itemize}
  \item \texttt{CallState.mode = listening} implies
    \texttt{CallAgentState.procState \in \{ready, dormant\}} -- no turn can be
    in flight while the call has not yet detected one to send.
  \item The transition into \texttt{CallState.mode = waiting}
    (\texttt{TurnDetected}) is exactly the caller-side trigger for
    \texttt{SendTurn} -- the two operations fire together, in the same
    caller-level step, though they are modelled separately because they
    belong to different objects with different lifetimes (\S\ref{sec:state}).
  \item \texttt{CallAgentState.procState = busy} therefore implies
    \texttt{CallState.mode \in \{waiting, speaking\}} -- once
    \texttt{ReplyBegins} fires on the audio side, the agent process may still
    be \texttt{busy} composing the remainder of its reply (FR-11: speech
    begins on the first complete portion, before the whole reply exists),
    so \texttt{speaking} does not require \texttt{ReplyComplete} to have
    fired yet.
  \item \texttt{CallState}'s own \texttt{EndCall} (hangup from any active
    mode, including mid-\texttt{speaking}) is exactly the caller-side trigger
    for \texttt{AbortTurn} when \texttt{CallAgentState.procState = busy}, or
    \texttt{Teardown} when it is \texttt{ready} -- \texttt{EndCall} alone
    does not determine which; the caller reads
    \texttt{CallAgentState.procState} first.
\end{itemize}

This correspondence is stated in prose, in the same manner
\texttt{conversation-mode-call-state.tex}'s \S``The Single Serialized
Dispatch Point'' states its own cross-cutting discipline in prose rather than
as a further schema: Z models the shape of state and the
preconditions/postconditions of a transition, not the scheduling discipline
between two collaborating state machines owned by the same caller.

\section{What This Model Does Not Cover}

\begin{itemize}
  \item \textbf{How read-only enforcement is implemented} -- a permission-mode
    flag, a sandboxed filesystem view, or harness-level tool restriction are
    all still open per DES-066; this model only fixes when the property must
    hold (\S\ref{sec:state}, invariant 2).
  \item \textbf{The context-snapshot's own content and construction} -- what
    goes into the opening message \texttt{Spawn} delivers is a
    \texttt{CallSession}-level concern (FR-4's task-context requirement),
    out of this process-lifecycle model's scope.
  \item \textbf{The post-call summary handoff to the primary session} -- DES-066
    lists this as undesigned; it happens after \texttt{Teardown} or
    \texttt{AbortTurn} reach \texttt{ended}, entirely outside this model.
  \item \textbf{Process-supervision mechanics} (a bare \texttt{subprocess.Popen}
    held by \texttt{voxd} itself, vs. \texttt{tmux}/\texttt{pi-tools}' \texttt{keep}
    extension) -- an implementation choice about \emph{how} \texttt{Spawn} and
    \texttt{Teardown} are realized, not a state or invariant this model needs
    to express.
\end{itemize}

\end{document}
